
Det antages at læseren har en grundlæggende forståelse for differentialligninger og numeriske metoder. 
Dette afsnit vil give en introduktion til det matematiske grundlag for PDE'er, samt de numeriske metoder der anvendes i BOUT++-simuleringer.
Derudover vil det også introducere moderne datadrevne metoder som PINNs og PINOs, som er relevante for dette projekt.
\cite{iter_new_energy}
\subsection{Det matematiske grundlag} \label{ch2_1}

% Primer til fysikken bag plasma modellen
Funktioner og funktionrum
    - Skalar- og vektorfunktioner
    - Funktioner af flere variable
    - Kontinuitet og differentiabilitet
    - Normer og afstandsmål
    - Gradient
    - Divergens
    - Rotation (curl)
    - Laplace-operatoren
    - Kædereglen for funktioner af flere variable

Klassifikation af PDE'er
    - Definition af PDE'er
    - Klassifikation af PDE'er
    - Orden og linearitet
    - Elliptiske PDE'er
    - Parabolske PDE'er
    - Hyperbolske PDE'er
    - Definer hvilket pde system vi har med at gøre

Begyndelses- og randbetingelser
    - Initialbetingelser (IC)
    - Dirichlet-betingelser
    - Neumann-betingelser
    - Periodiske randbetingelser

% Primer til simulering uden BOUT++ fokues
Numerisk approksimation
    - Diskretisering af kontinuert rum
    - Finite Difference approksimationer
    - Fremad differens
    - Bagud differens
    - Central differens
    - Diskretisering af tidsafledte
    - Diskret repræsentation af PDE'er

Fejl og konvergens
    - Trunkeringsfejl
    - Konsistens
    - Stabilitet
    - Konvergens

% Primer til PINN/PINO-fokus
-   Funktioner og funktionrum.
-   Partielle afledte og differentialoperatorer.
-   PDE'er og randbetingelser.
-   Finite Difference-approksimationer.
-   Automatisk differentiation.
-   Residualer og MSE-tab.

% udledning af PDE-systemet
\subsection{Fysiken bag Hesel modellen}
* Simulering af fysiske systemer
  * Bevarelseslove
    * Masse
    * Momentum
    * Energi
  * Advektion, diffusion og reaktion
  * Turbulens og kaotisk dynamik
  * Multiskala-problemer
  * Plasmafysiske PDE-systemer

% numeriske metoder til løsning/simulering af PDE'er
% og hvordan BOUT++ og BOUT-HESEL implementerer disse og hvordan det faktsik fungere på en computer
\subsection{Numeriske metoder, BOUT-HESEL simuleringer}
% Viderebyg på det eksisterende afsnit om det matematiske grundlag for PDE'er, men nu med fokus på BOUT++-simuleringer 
% Du skal også begynde at inkludere computer science elementer

* Numeriske metoder til løsning/simulering af PDE'er

  * Diskretisering af rum og tid

    * Uniforme og ikke-uniforme grids
    * Stabilitet og konvergens
  * Finite Difference Method (FDM)

    * Fremad-, bagud- og centraldifferenser
    * Numerisk approksimation af afledte
  * Finite Volume Method (FVM)

    * Bevarelseslove
    * Flux-beregninger
  * Finite Element Method (FEM)

    * Svag formulering
    * Basisfunktioner
  * Spektrale metoder

    * Fourier-serier
    * Fourier-transformer
    * Pseudospektrale metoder
  * Tidsintegration

    * Euler-metoder
    * Runge-Kutta-metoder
    * Implicitte og eksplicitte metoder
  * Stabilitetsanalyse
    * CFL-betingelsen
    * Fejl- og konvergensanalyse

% SciML, PINN og PINO

\subsection{Datadrevne og fysik informeede metoder}

* Moderne datadrevne metoder

  * Physics-Informed Neural Networks (PINNs)

    * Automatisk differentiation
    * PDE-residualer
    * IC- og BC-tab
  * Neural Operators

    * Fourier Neural Operators (FNO)
    * Operatorlæring
    * Generalisering mellem diskretiseringer
  * Hybridmodeller

    * Kombination af simulering og data
    * Data-assimilering

% Hardconstraints og soft constraints i PINNs


% Måske også en gennemgang af computer hardware